
\subsubsection{6.1 Key Findings}

\textbf{Finding 1: Style-function dissociation as the primary mechanism.} Total coverage without category analysis is misleading for evaluating feedback quality. Pylint's 100\% total coverage of HumanEval failures overstates its diagnostic value by concealing that 94.3\% of flags are Convention-category style rules. When the LLM receives ''fix your formatting'' as its primary repair guidance for an algorithmic failure, it may restructure code while addressing style, introducing new logical errors in the process. This mechanism explains the observed HumanEval regression under pylint repair. Future work reporting pylint coverage of LLM failures should decompose by flag category to distinguish functional from stylistic coverage.

\textbf{Finding 2: Execution feedback is practically effective at fixed compute.} The +40.2pp MBPP improvement from a single round at B=1000 output tokens is a substantively large effect --- the baseline success rate more than doubles on a standard benchmark at minimal inference cost. The +4.9pp HumanEval improvement is smaller in magnitude but statistically robust (McNemar p=0.0001). Execution feedback is a reliable choice for single-round repair across both benchmark types examined here.

\textbf{Finding 3: Task complexity moderates static analysis feedback utility.} Even feedback with low functional informativeness (94.3\% style flags) produces positive gains on MBPP's simple function-completion tasks (+18.3pp). The observation that style-guided rewrites preserve logical structure in short functions but disrupt it in complex algorithms suggests that the appropriate feedback signal may depend on the complexity profile of target tasks.

\subsubsection{6.2 Limitations}

\textbf{L1: Single model.} All results are based on Llama 3.1 8B Instruct only. A planned Qwen2.5-Coder-7B replication experiment (h-m3) was not executed due to resource constraints (single GPU availability during the h-m1 execution phase). Claims are scoped to Llama 3.1 8B Instruct; generalizability to other 7B-scale models, including code-specialized variants, is not confirmed.

\textbf{L2: Single effective repair round.} At B=1000, at most one repair round completes for most problems. Results characterize single-round repair, not multi-round iterative improvement. Per-round trajectory data confirms that rounds 2 and 3 contribute near-zero incremental improvement, suggesting that B=1000 captures most of the available single-model improvement, but multi-round dynamics at larger budgets (B=2000+) are not examined.

\textbf{L3: Primary prediction for coverage refuted (informative null).} The original prediction for h-m2 was that pylint/mypy would flag fewer than 50\% of HumanEval baseline failures. Actual total coverage is 100\%. This null result on the primary prediction metric is reported transparently. The functional coverage result (12.5\% E+W) is a separate, positive finding that supports the same underlying mechanistic argument and is arguably more precise than the original prediction.

\textbf{L4: Mechanism is inferred from aggregate results.} The style-guided corruption explanation for the HumanEval regression ($\Delta _pylint_HE = -4.3pp)$ is consistent with the aggregate data but is not confirmed by per-problem analysis. A per-problem comparison of round 0 and round 1 solutions for the 7 HumanEval problems that regressed under pylint repair would be needed to confirm whether structural rewrites correlate with regression.

\subsubsection{6.3 Broader Impact}

This work contributes empirical evidence relevant to the responsible deployment of LLM-based code generation tools. Production systems using static analysis output as repair feedback should be evaluated on functional correctness metrics, not total coverage. The style-function dissociation documented here may be invisible in aggregate quality metrics but visible in pass@k evaluations. No significant negative societal impacts are identified from this benchmarking study.

